\documentclass[11pt,a4paper,notitlepage]{report}
% Motion-gated dynamic-object masking: implementation and results.
%
% Standalone technical report on `filter_mode:=motion` in ORB-SLAM3 and
% VINS-Fusion.  Format follows the final report's conventions (Times 12 pt
% body, booktabs tables, siunitx units) so that passages can be lifted into
% docs/report/final_report/chapters/ without reformatting, but it compiles on
% its own:   latexmk -pdf filter_motion.tex
%
% Source of every number is cited in the text to a run directory under
% output/ or to an engineering-record section.  Numbers marked "pooled" are
% regenerated by script/analyse_motion_gate.py into
% output/compare/simulation/exp3-motion-gate/.

\usepackage[T1]{fontenc}
\usepackage[utf8]{inputenc}
\usepackage{mathptmx}
\usepackage[scaled=0.92]{helvet}
\usepackage{courier}
\usepackage[a4paper,left=30mm,right=20mm,top=25mm,bottom=25mm]{geometry}
\usepackage{microtype}
\usepackage{booktabs}
\usepackage{tabularx}
\usepackage{array}
\usepackage{amsmath}
\usepackage{siunitx}
\usepackage{graphicx}
\usepackage{caption}
\usepackage{enumitem}
\usepackage[table]{xcolor}
\PassOptionsToPackage{hyphens}{url}\usepackage{url}
\usepackage[hidelinks]{hyperref}
\usepackage[nameinlink,noabbrev]{cleveref}

\newcommand{\repo}[1]{\texttt{\detokenize{#1}}}
\sisetup{detect-all,group-separator={,}}
\DeclareSIUnit{\pixel}{px}
\captionsetup{font=small,labelfont=bf}
\setlength{\parskip}{4pt}

\title{Motion-Gated Dynamic-Object Masking\\
       \large Implementation and Evaluation in ORB-SLAM3 and VINS-Fusion}
\author{Bolton Athit Davies}
\date{Reporting cutoff: 21 October 2026}

\begin{document}
\maketitle

\begin{abstract}
\noindent
Class-based masking removes every detected bin, box, bucket and pallet jack
from the visual front end, whether or not the object is moving. In the
warehouse datasets used here, most detected objects are stationary scenery, so
the filter discards good landmarks to remove a minority of harmful ones. This
report documents \repo{filter_mode:=motion}, an extension that masks a detected
box only after the tracker has observed it move, its implementation in both
pipelines, and its evaluation against simulator ground truth.

\textbf{The result is negative.} In ORB-SLAM3 the gate is at chance: pooled over
\num{7710} ground-truth-labelled box-frames it flags \SI{82}{\percent} of
genuinely moving objects and \SI{76}{\percent} of stationary ones, an area
under the ROC curve of \num{0.528} where \num{0.5} is chance. In VINS-Fusion
three successive signals were tried and all failed, one of them inverted.
Both failures share a cause that is architectural rather than a matter of
tuning: each estimator discards the features that disagree with its own pose
estimate, so by the time a census can be taken, the evidence of motion has
already been removed. The feature is implemented, documented, disabled by
default, and reported here as a measured negative result.
\end{abstract}

\tableofcontents

\chapter{Motivation}

\section{What class masking does}

The dynamic-object filter implemented for Experiment~3 follows RY-SLAM: a YOLO
detector publishes bounding boxes, and the SLAM front end refuses to use image
features that fall inside them. In ORB-SLAM3 the mask is consumed inside
\repo{ORBextractor}, which discards keypoints landing on a dynamic region
before the octree distributes the per-level feature quota. In VINS-Fusion it is
consumed inside \repo{FeatureTracker::setMask()}.

The filter is \emph{semantic}: it masks objects whose \emph{class} is one the
detector was trained on. Class membership is not evidence of motion. A pallet
jack parked against a wall is masked exactly like one being driven across the
aisle.

\section{Why that matters in these datasets}

The proportion of detected objects that actually move is small. Measured from
\repo{/world/default/pose/info} by first-versus-last displacement
(\SI{>5}{\centi\meter}, excluding the robot):

\begin{table}[htbp]
  \centering
  \caption{Moving props per dataset, measured from simulator ground truth.}
  \label{tab:movers}
  \small
  \begin{tabular}{llrr}
    \toprule
    Dataset & World & Entities & Moving \\
    \midrule
    \repo{dataset_allsensor_realistic_000} & U static twin        & 131 & 0 \\
    \repo{dataset_allsensor_realistic_001} & I static twin        & 153 & 0 \\
    \repo{dataset_allsensor_realistic_002} & U dynamic            & --- & 10 \\
    \repo{dataset_allsensor_realistic_003} & I dynamic            & --- & 12 \\
    \repo{dataset_dynamic_nofloortexture_00_000} & dynamic \num{00} & 101 & 25 \\
    \bottomrule
  \end{tabular}
\end{table}

In the realistic worlds, ten to twelve props move among roughly \num{150}
static ones. Class masking therefore spends most of its effect deleting
landmarks that were never a problem. In the static twins it can only do harm:
there is nothing to remove.

This is the gap motion gating is meant to close --- keep the parked objects as
landmarks, mask only the ones in motion.

\section{Relation to IQ-VIO}

IQ-VIO \cite{iqvio} addresses the same problem from the opposite direction: it
keeps every feature and instead lowers the confidence of the visual pose
measurement inside an EKF when the coverage of dynamic-class detections is
high. That protects the fused output but not the map, and --- importantly ---
its interference index is computed from detection \emph{coverage}, so like
class masking it never measures motion. A parked forklift and a moving one
produce the same index. The approach reported here differs in deciding
per object, and in using the tracker's own geometry rather than detection area.

\chapter{Implementation}

\section{Where the mask is consumed, and why it differs per pipeline}

The two estimators manage features differently, and the integration point
follows from that difference rather than from preference.

\paragraph{ORB-SLAM3 --- one population.}
Features are re-extracted from scratch every frame, so masking extraction is
the whole job. The mask is consumed inside \repo{ORBextractor}, in the loop
that collects candidate keypoints before the octree distributes the per-level
quota. The rule is RY-SLAM equation~(1): count dynamic pixels in a
$3\times3$ neighbourhood of the keypoint and delete it at five or more,
reserve it otherwise. One threshold covers both cases the paper separates ---
a keypoint deep inside a dynamic region scores nine and goes, one just outside
a boundary scores at most four and stays.

Two coordinate subtleties make this less trivial than it reads. At that point
in the extractor a keypoint's position is in \emph{level} coordinates relative
to \repo{(minBorderX, minBorderY)}, not image coordinates, so the offset must
be added back before indexing the mask. And the neighbourhood is spaced by the
level's scale factor rather than by one pixel, so a coarse-pyramid keypoint ---
whose true support in the source image is larger --- is tested over a
correspondingly larger area:

With $s$ the level's scale factor and $(x_0,y_0)$ the keypoint mapped back into
image coordinates,
\[
  x_0 = (\mathit{pt}.x + \mathit{minBorderX})\,s, \qquad
  y_0 = (\mathit{pt}.y + \mathit{minBorderY})\,s,
\]
\[
  n_{\text{dyn}} = \bigl|\;\{(dx,dy) \in \{-1,0,1\}^2 \;:\;
      \mathrm{mask}\left[y_0 + dy\,s,\; x_0 + dx\,s\right] = 0 \}\;\bigr|,
\]
and the keypoint is discarded when $n_{\text{dyn}} \ge 5$. Samples falling
outside the image are skipped rather than treated as dynamic.

An empty mask short-circuits the test entirely, so a run with filtering off is
the stock tracker bit for bit and the recorded baseline stays reproducible.

\paragraph{VINS-Fusion --- two populations.}
Features come from two places. \emph{Tracked} points are propagated by
Lucas--Kanade optical flow and are never re-detected; a point that latches onto
a moving crate rides it for hundreds of frames. \emph{New} points top the set
back up to \repo{max_cnt} wherever there is room. Masking only new detections
would leave the first population untouched and the filter nearly useless.

\repo{FeatureTracker::setMask()} sits exactly between the two. Upstream already
builds a canvas there, walks the tracked points in descending track length,
keeps those whose pixel reads \num{255}, paints a \repo{min_dist} exclusion
circle at each kept point, and then hands that same \repo{cv::Mat} to
\repo{goodFeaturesToTrack} as its mask argument. Seeding the canvas from the
dynamic mask instead of a blank all-\num{255} image therefore does both jobs
with one substitution: a tracked point inside a box fails the existing test and
is dropped from \repo{cur_pts/ids/track_cnt}, and a new point cannot be seeded
there either. No new loop and no index bookkeeping ---
\repo{prev_pts} is rebuilt wholesale from \repo{cur_pts} at the end of
\repo{trackImage}, so dropping extra points breaks no correspondence.

The mask is applied to \repo{image0} only. The right image's points are tracked
\emph{from} the already-masked left set, so nothing re-enters through the
second camera.

\section{Mask construction}

Detections arrive as a \repo{vision_msgs/Detection2DArray} stamped with the
\emph{source image's} timestamp, never \repo{now()}; using publication time
would offset every mask by the inference latency and drag the masks behind the
objects. The node buffers them by that stamp and, for each frame, selects the
detection nearest in time within \repo{det_max_age}.

The mask itself is \repo{CV_8UC1}, \num{255} meaning keep. Each box is inflated
by \repo{mask_dilate_px} (\num{8}) before painting, because a feature on the
silhouette of a moving crate is exactly the one that corrupts the map, and
detector boxes clip their objects slightly. Covered area is counted
\emph{before} painting so that overlapping boxes are not double-counted.

A safety valve discards the whole mask when it would cover more than
\repo{max_mask_fraction} (\num{0.8}) of the frame: one runaway false positive
spanning the image would starve the tracker and lose the map, which is far
worse than not filtering that frame. The threshold is measured rather than
guessed --- over all \num{791} frames of an earlier dynamic bag the dilated
boxes covered a mean of \SI{43.5}{\percent}, p90 \SI{61}{\percent}, maximum
\SI{72.8}{\percent}, so a limit of \num{0.6} would have fired during entirely
normal operation.

When a frame's detection contains no boxes the node returns an \emph{empty}
\repo{cv::Mat} rather than an all-\num{255} one. That is both correct and
cheaper: the extractor short-circuits on it instead of running the
per-keypoint test.

\section{The census}

Motion gating adds a read-back path that the class filter does not need.

\paragraph{ORB-SLAM3.} The fork's \repo{System::TrackStereo} gains an optional
\repo{voteBoxes} argument and a \repo{GetBoxVotes()} accessor. After tracking
returns, every left keypoint of the current frame is assigned to the first box
containing it, or to the background if it lies outside all of them, and two
counters are incremented: keypoints, and keypoints that ended the frame as
inliers.

The boxes are given in original image coordinates and must be scaled by the
same factor the images were, since \repo{TrackStereo} may rectify or resize
before feeding the tracker. Overlapping boxes credit the keypoint to the first
match, which keeps the partition exact --- every keypoint is counted once, so
the background count is a true complement.

``Inlier'' is deliberately narrow. A map point qualifies only when it

\begin{enumerate}[leftmargin=*,itemsep=1pt]
  \item is attached to the keypoint and not flagged an outlier by pose
        optimisation;
  \item satisfies \repo{Observations() > 0}, i.e.\ is anchored at a keyframe.
        ORB creates \emph{temporary} map points from the previous frame's
        stereo depth on every frame; at the movers' speed those match within
        the projection search window and would mask the very motion being
        looked for;
  \item is at least \repo{motion_min_point_age} frames old
        (\repo{System::SetBoxVoteMinAge}), because local mapping inserts a
        keyframe roughly every ten frames and re-anchors fresh points on a
        mover, so ``keyframe-anchored'' alone is only about
        \SI{0.3}{\second} of baseline.
\end{enumerate}

\paragraph{VINS-Fusion.} There is no per-frame inlier flag, so the analogue is
whether the sliding window holds a solved, positive depth for the feature
(\repo{FeaturePerId::isSolved()}). The census runs on the backend thread
immediately after \repo{processImage()} --- the only instant at which the
window's feature list and this frame's observations are both current --- and
pushes its result onto a bounded queue that the node drains at the start of the
next \repo{maskFor()}.

\section{Concurrency and the in-flight ledger}

The verdict for a frame is produced on a different thread from the one that
consumed its mask, so the two must be reconciled by timestamp.

In ORB-SLAM3 the mask is built in the ROS subscription callback, the frame is
pushed to a single-slot queue, tracking runs on the wrapper's worker thread,
and the census returns through the result callback. In VINS the mask is built
in the image callback and the census is produced by the backend thread. In both
cases the node keeps an \emph{in-flight ledger}: a map from frame timestamp to
the boxes handed out for that frame, together with the verdict each inherited.
When a census arrives, its timestamp is looked up, the verdicts are updated,
and every older entry is discarded --- those frames will never be scored.

Two consequences follow, and both are properties of the design rather than
bugs:

\begin{itemize}[leftmargin=*,itemsep=2pt]
  \item The ledger is bounded (\num{30} entries). If the census falls further
        behind than that, entries are evicted unscored and their boxes remain
        \emph{unknown} permanently.
  \item A frame dropped before tracking never produces a census, so its ledger
        entry is discarded when the next census overtakes it.
\end{itemize}

All ledger access is under one mutex, and the census carries raw counts only
--- no policy crosses the thread boundary, which keeps the decision rule
entirely inside the node and changeable without touching either estimator.

\section{Verdict, association and probation}

With $r = n_{\text{inlier}}/n_{\text{kp}}$ for a box and $r_{\text{bg}}$ for
the background of the same frame, the box is judged \emph{moving} when
$r < \alpha\,r_{\text{bg}}$ and \emph{static} otherwise. The comparison is
against the background of the \emph{same} frame rather than an absolute
threshold, because the inlier rate varies with viewpoint, texture and tracking
health; the background is the per-frame control.

The vote is withheld, and the previous verdict inherited, when the box holds
fewer than \repo{motion_min_points} keypoints or when its association chain is
younger than \repo{motion_min_box_age}. The second gate exists because a
one-frame census cannot distinguish \emph{moving} from \emph{newly revealed}:
both have only young map points. Requiring a second of continuous association
lets a static object accumulate old points before it is judged.

Each box of the next frame inherits the verdict of the previous frame's box
with the highest intersection over union, provided it exceeds
\repo{motion_iou}; otherwise it is \emph{unknown} and, by default, left
unpainted. Only \emph{moving} boxes are painted, so the filter is conservative:
anything it has not positively convicted keeps its features.

Probation closes a loop that would otherwise be one-way. A painted box contains
no keypoints, so its verdict could never be revisited and a single false
positive would propagate through association indefinitely --- which is exactly
what happened in the first run (\cref{sec:run1}). Every
\repo{motion_revote_frames} frames, a chain that is currently masked is left
unpainted for one frame so that it is scored again. The cost is one frame of
leakage per interval for a true mover; the benefit is that no verdict is
permanent.

\section{Parameters}

\begin{table}[htbp]
  \centering
  \caption{Motion-gate parameters and their calibrated defaults. Identical in
           both pipelines except \repo{motion_min_point_age}, which has no
           VINS analogue.}
  \label{tab:params}
  \small
  \begin{tabularx}{\textwidth}{llX}
    \toprule
    Parameter & Default & Meaning \\
    \midrule
    \repo{filter_mode}            & \repo{class} & \repo{class} masks every detected box; \repo{motion} masks only boxes seen to move \\
    \repo{motion_iou}             & \num{0.3}  & IoU required to associate a box with the previous frame's verdict \\
    \repo{motion_min_points}      & \num{15}   & fewer keypoints in the box than this and the vote is withheld \\
    \repo{motion_inlier_ratio}    & \num{0.5}  & $\alpha$: moving when the box's inlier ratio falls below this multiple of the background's \\
    \repo{motion_min_box_age}     & \num{30}   & frames a box must be associated before it is judged \\
    \repo{motion_min_point_age}   & \num{30}   & ORB only: map points younger than this do not count as census inliers \\
    \repo{motion_revote_frames}   & \num{15}   & frames between probation openings of a masked box \\
    \repo{motion_unknown_mask}    & false      & whether a never-judged box is masked \\
    \midrule
    \repo{mask_dilate_px}         & \num{8}    & box inflation before painting \\
    \repo{det_max_age}            & \SI{0.15}{\second} & oldest detection accepted for a frame \\
    \repo{max_mask_fraction}      & \num{0.8}  & discard the mask above this frame coverage \\
    \bottomrule
  \end{tabularx}
\end{table}

\section{Instrumentation}

Every quantity the evaluation needs is written by the run itself, so no number
is inferred from unrelated logs.

\begin{itemize}[leftmargin=*,itemsep=2pt]
  \item \repo{motion_votes.csv} --- one row per box per scored frame: box
        corners, the raw census (\repo{n_kp}, \repo{n_inlier} and the
        background pair), the inherited verdict, the final verdict and the
        chain age. This is the file the ground-truth scoring consumes.
  \item \repo{yolo_mask.csv} --- one row per frame: number of boxes painted,
        mask coverage, whether a detection matched, its age, whether the mask
        was rejected by the valve, and the \repo{n_moving}/\repo{n_static}/
        \repo{n_unknown} split. Rows are flushed individually, because a run is
        normally ended with a signal that does not run destructors and an
        unflushed stream loses its tail.
  \item \repo{run_summary.csv} / \repo{filter_summary.csv} ---
        \repo{filter_mode}, voted totals, \repo{detection_match_rate} and
        \repo{filter_status} (\repo{none}, \repo{never_fired}, \repo{degraded},
        \repo{ok}).
  \item VINS additionally separates \repo{dynamic_points_dropped} from
        \repo{mindistance_points_dropped}. Both kinds of rejection happen in
        the same loop and an earlier version attributed all of them to the
        filter, which overstated its effect; the split tests the original mask
        at each rejected point and is exact.
\end{itemize}

\section{Evaluation method}

\repo{script/motion_votes_eval.py} scores one run. For each row of
\repo{motion_votes.csv} it reconstructs the camera pose at that timestamp,
projects every prop's collision mesh into \repo{cam0}, takes the 2-D extent as
the truth box, matches the voted box to a prop by IoU, and labels that prop
moving or static from its world-frame speed over a short window. Two traps are
handled explicitly: the simulator publishes \repo{/world/default/pose/info}
with all-zero header stamps, so simulation time is rebuilt by interpolating
through \repo{/clock}; and the extrinsic matrix is parsed by hand because
\repo{cv::FileStorage} mis-reads the particular block in these configuration
files.

\repo{script/analyse_motion_gate.py} answers the prior question --- whether a
usable signal exists at all --- by pooling every labelled box-frame across runs
and sweeping both gates. It reports the area under the ROC curve and Youden's
$J$ rather than $F_1$, because roughly three quarters of labelled box-frames
are moving, so a rule that calls everything moving scores $F_1 \approx 0.86$
while being useless. Its outputs are written to
\repo{output/compare/simulation/exp3-motion-gate/}.

\chapter{Results}

\section{Validity of the runs scored here}

Two classes of defect were found during this work and are stated before the
results, because they bound what the numbers can support.

\textbf{Mask availability varies with machine load.} The detector runs as a
separate process competing with the tracker. On this hardware the fraction of
frames that received a mask ranged from \SI{1}{\percent} to
\SI{100}{\percent} across repeats of one bag, declining within a session as the
laptop heated. Every filtered result must therefore be read beside
\repo{detection_match_rate} and \repo{filter_status}: a low match rate means the
filter did not run, not that filtering did not help. The two runs pooled in
\cref{sec:orbneg} both carry \repo{filter_status: degraded}
(match rate \num{0.78} and \num{0.87}).

\textbf{A census cannot be compared across decision rules.} It is computed
downstream of its own masking, so changing the rule changes which features
survive and therefore changes the census. Numbers from runs that used different
verdict rules are not comparable, however similar their configuration.

\section{ORB-SLAM3: four design iterations}

The gate was developed over four runs on \repo{dataset_allsensor_realistic_003}
(\num{12} moving props), each fixing the fault the previous one exposed.

\subsection{Run 1 --- runaway}
\label{sec:run1}

With no probation, a single false \emph{moving} verdict propagated through IoU
association until motion mode had collapsed into class mode at
$t \approx \SI{105}{\second}$: \SI{81}{\percent} of static box-frames were
masked. Separately, the census counted ORB's \emph{temporary} map points, which
are created from the previous frame's stereo depth; at the movers' speed of
\SIrange{0.17}{0.27}{\meter\per\second} the apparent motion is about
\SI{2}{\pixel} per frame, well inside the projection search window, so a slow
mover's temporary points matched and counted as inliers.

\subsection{Runs 2--4 --- the progression}

\begin{table}[htbp]
  \centering
  \caption{ORB-SLAM3 motion gate on \repo{dataset_allsensor_realistic_003}.
           Per-vote figures are measured only on box-frames that reached a
           vote; the masked columns describe what the tracker actually
           experienced.}
  \label{tab:orbruns}
  \small
  \begin{tabular}{llrrrr}
    \toprule
    Run & Change & \multicolumn{2}{c}{Masked box-frames} & \multicolumn{2}{c}{Per vote} \\
    \cmidrule(lr){3-4}\cmidrule(lr){5-6}
        &        & moving & static & recall & false alarm \\
    \midrule
    1 & defaults                       & ---  & \SI{81}{\percent} & 0.39 & \SI{7.7}{\percent} \\
    2 & keyframe-anchored + probation  & ---  & \SI{53}{\percent} & 0.49 & \SI{9.0}{\percent} \\
    3 & + map-point age floor          & ---  & \SI{62}{\percent} & 0.62 & \SI{15.4}{\percent} \\
    4 & + box chain age (defaults)     & \SI{56}{\percent} & \SI{30}{\percent} & 0.43 & \SI{5.0}{\percent} \\
    \bottomrule
  \end{tabular}
\end{table}

Run~3 is the instructive one. The map-point age floor did what it was designed
to do on movers --- their median inlier ratio fell to a clean \num{0.00} ---
and did exactly the same to \emph{static objects that had just come into view},
whose map points are equally young. The false-alarm rate doubled. The
one-frame census cannot distinguish \emph{moving} from \emph{new}; the confound
is scene novelty. Run~4's chain-age gate removes that confound by refusing to
judge a box until it has been in view for a second, at the cost of the recall
the age floor had bought.

At this point the gate masked \SI{56}{\percent} of moving box-frames while
sparing \SI{70}{\percent} of the static landmarks that class mode destroys.
That appeared to be a usable, if mediocre, trade.

\section{ORB-SLAM3: the result does not survive a larger sample}
\label{sec:orbneg}

The runs above rest on small per-vote samples --- run~4's false-alarm figure
comes from the subset of box-frames that passed both gates. To test whether the
trade was real, the gate was scored on
\repo{dataset_dynamic_nofloortexture_00_000}, which carries \num{25} moving
props, pooling both available runs and sweeping both gates so that the static
class is large enough to estimate a false-alarm rate from.

Pooled over \num{7710} ground-truth-labelled box-frames
(\num{5833} moving, \num{1877} static), what the mask was built from:

\begin{table}[htbp]
  \centering
  \caption{Pooled verdicts at the shipped thresholds,
           \repo{dataset_dynamic_nofloortexture_00_000}. Source:
           \repo{output/compare/simulation/exp3-motion-gate/}.}
  \label{tab:pooledconf}
  \small
  \begin{tabular}{lrrrr}
    \toprule
    Truth & unknown & static & moving & masked \\
    \midrule
    moving & \num{4928} & \num{116} & \num{789} & \SI{13.5}{\percent} \\
    static & \num{1561} & \num{14}  & \num{302} & \SI{16.1}{\percent} \\
    \bottomrule
  \end{tabular}
\end{table}

The gate masks stationary objects slightly \emph{more} often than moving ones.
Sweeping both gates over the same pooled data shows this is not a threshold
choice:

\begin{table}[htbp]
  \centering
  \caption{Separability of the ORB census, pooled.
           $J = \text{recall} - \text{false alarm}$; AUC is
           $P(\text{moving ranks below static})$, so \num{0.5} is chance.
           \repo{motion_min_points} held at \num{15}.}
  \label{tab:sweep}
  \small
  \begin{tabular}{rrrrrrr}
    \toprule
    \repo{min_box_age} & coverage & $n_{\text{mov}}$ & $n_{\text{stat}}$ & recall & false alarm & AUC \\
    \midrule
    \num{0}  & \SI{56.4}{\percent} & \num{3307} & \num{1044} & 0.820 & 0.760 & 0.528 \\
    \num{3}  & \SI{41.5}{\percent} & \num{2483} & \num{714}  & 0.809 & 0.770 & 0.519 \\
    \num{5}  & \SI{34.6}{\percent} & \num{2088} & \num{582}  & 0.793 & 0.790 & 0.506 \\
    \num{10} & \SI{23.2}{\percent} & \num{1398} & \num{390}  & 0.737 & 0.821 & 0.457 \\
    \num{15} & \SI{17.0}{\percent} & \num{994}  & \num{316}  & 0.682 & 0.804 & 0.423 \\
    \num{20} & \SI{11.9}{\percent} & \num{664}  & \num{250}  & 0.640 & 0.784 & 0.400 \\
    \num{30} & \SI{3.0}{\percent}  & \num{188}  & \num{43}   & 0.394 & 0.605 & 0.437 \\
    \bottomrule
  \end{tabular}
\end{table}

Recall and false alarm move together at every setting. The best separation
anywhere in the sweep is $J = +0.061$ --- recall \num{0.820} at
\SI{76.0}{\percent} false alarm. Loosening \repo{motion_min_box_age} raises
coverage from \SI{3}{\percent} to \SI{56}{\percent}, confirming that the chain
age, not the feature count, is what starves the gate; but the extra samples do
not improve separation, they dissolve it. At the shipped setting the AUC is
\num{0.437}, slightly \emph{below} chance.

\section{VINS-Fusion: three signals, all negative}

The same policy on \repo{dataset_dynamic_nofloortexture_00_000}, where
\SI{64}{\percent} of box-frames sit on genuinely moving objects:

\begin{table}[htbp]
  \centering
  \caption{Census signals tried in VINS-Fusion and their outcomes.}
  \label{tab:vins}
  \small
  \begin{tabularx}{\textwidth}{lXl}
    \toprule
    Signal & Quantity & Outcome \\
    \midrule
    solved depth & feature holds a solved, positive depth in the window
      & recall \num{0.25} at \SI{23}{\percent} false alarm \\
    reprojection residual & median per-box residual against the optimised window
      & \textbf{inverted}: movers score lower \\
    residual availability & fraction of a box's features that have a residual
      & recall \num{0.56} at \SI{41}{\percent} false alarm \\
    \bottomrule
  \end{tabularx}
\end{table}

Feature starvation was ruled out explicitly. Raising \repo{max_cnt} from
\num{150} to \num{250} had no effect --- the scene cannot supply more corners
at the configured spacing, and the per-box feature distribution was unchanged.
Reducing \repo{min_dist} from \num{30} to \SI{15}{\pixel} tripled the features
per box, raising coverage from \SI{19.6}{\percent} to \SI{55.8}{\percent}, and
the signal \emph{vanished}: \SI{94.7}{\percent} of features inside moving
objects held solved depths, and not one of \num{83} voted rows fell below the
threshold. More data made the result worse, which is the signature of a
quantity that was never measuring the intended property.

\chapter{Why both pipelines fail}

The three VINS signals and the ORB census share one structural property: all of
them are read from estimator state that has already been cleaned of moving
points.

\begin{itemize}[leftmargin=*,itemsep=2pt]
  \item ORB-SLAM3's census counts map points that survived \repo{TrackLocalMap}.
        Points inconsistent with the optimised pose have already been marked
        outliers. What remains inside a moving object's box is whatever happened
        to fit.
  \item VINS-Fusion's \repo{outliersRejection} purges any feature whose average
        reprojection error exceeds \SI{3}{\pixel}, every frame. A mover's
        disagreeing features are gone within a frame or two, so the residual
        census measures survivors --- which is why that signal came out
        inverted.
  \item A moving point also triangulates perfectly well. It converges to the
        \emph{wrong} depth, but a wrong depth is still a solved depth, and no
        test on depth validity can separate the two.
\end{itemize}

Both estimators are designed to be robust to moving points. That robustness is
exactly what makes them poor instruments for detecting motion: by the time
either has an opinion about a feature, the features that disagreed with it have
been discarded. This is not a tuning failure and no threshold on these
quantities can repair it.

A signal that could work must be computed \emph{before} rejection. Three
candidates exist in the VINS front end and are presently unused:
\repo{rejectWithF()}, a RANSAC fundamental-matrix check that is written but
commented out, whose outliers are by definition points violating epipolar
geometry; \repo{predictPtsInNextFrame()}, implemented but never called, which
would give a predicted-versus-tracked displacement per feature; and
\repo{ptsVelocity}, a per-point image velocity already computed every frame, in
which ego-motion is a smooth field and a mover a coherent local anomaly. The
ORB-SLAM3 equivalent is the reprojection residual of each keypoint against the
IMU-predicted pose, computed in \repo{Tracking} before outlier marking. None of
these has been tried.

\chapter{Status and recommendation}

\textbf{Status: implemented, evaluated, negative.} The feature is complete and
in the tree in both pipelines, disabled by default
(\repo{filter_mode: class}), with calibration tooling and ground-truth scoring.
It does not detect motion in either pipeline.

Experiment~3 should therefore be reported on \repo{class} mode, which is the
RY-SLAM method the project set out to evaluate. Motion gating belongs in the
conclusions as an attempted extension with a measured negative result and an
identified architectural cause.

Two cautions for anyone extending this work. First, every encouraging result in
this project's motion-gate history came from a small sample and did not survive
a larger one: a separation measured on \num{19} static rows disappeared at
\num{1044}, and a VINS recall of \num{0.62} rested on \num{16} moving rows.
No threshold or rate should be quoted before \repo{motion_votes_eval.py} or
\repo{analyse_motion_gate.py} has scored it against ground truth with the
sample size stated. Second, the per-vote figures in \cref{tab:orbruns} describe
only the box-frames that reached a vote; the masked columns, not the per-vote
columns, describe what the tracker experienced.

\begin{thebibliography}{9}
\bibitem{iqvio}
Zhang, H., Ye, F., Lai, Y., Li, K., and Xu, J., 2023,
``IQ-VIO: Adaptive Visual Inertial Odometry via Interference Quantization
under Dynamic Environments,''
\textbf{Intelligent Service Robotics}, Vol.~16, pp.~565--581.
\bibitem{ryslam}
Chen, Y., et al., 2025,
``RY-SLAM: A Robust YOLO-Based Real-Time Visual SLAM Solution for Dynamic
Environments,''
\textbf{Proceedings of ICVRV}.
\end{thebibliography}

\end{document}
